% Vietnamese translation for OI Public Library
% Source-PDF: IOI中国国家候选队论文/2002/王知昆.pdf
\documentclass[11pt]{article}
\usepackage{oipl-vn}

\title{Lựa chọn thứ tự tìm kiếm}
\author{Wang Zhikun}
\date{}

\begin{document}
\maketitle

\origtitle{Choosing the search order}
\sourcepath{IOI China candidate team papers / 2002 / Wang Zhikun.pdf}

\section{Mở đầu: bài toán xếp cách quãng}

Xét bài toán \term{xếp cách quãng}. Có \(2N\) khối gỗ; mỗi khối được ghi một
số tự nhiên từ \(1\) đến \(N\), và mỗi số xuất hiện trên đúng hai khối. Ta
cần xếp các khối thành một hàng sao cho giữa hai khối mang số \(i\) có đúng
\(i\) khối khác.

Chẳng hạn với \(N=3\), một cách xếp hợp lệ là
\[
  3,\ 1,\ 2,\ 1,\ 3,\ 2.
\]
Bài toán lập trình là: với \(N\le 40\) cho trước, hãy tìm một cách xếp hợp lệ.

Chỉ thay đổi thứ tự tìm kiếm đã tạo ra khác biệt rất lớn về thời gian chạy:

\begin{center}
\begin{tabular}{c|cc}
\hline
\(N\) & Tìm từ lớn đến nhỏ & Tìm từ nhỏ đến lớn \\
\hline
11 & 0.00 sec & 0.00 sec \\
15 & 0.00 sec & 0.11 sec \\
20 & 0.00 sec & 36.32 sec \\
32 & 0.00 sec & quá thời gian \\
40 & 0.00 sec & quá thời gian \\
\hline
\end{tabular}
\end{center}

Ví dụ này cho thấy thứ tự xét các phần tử tìm kiếm không phải là chi tiết phụ.
Nó có thể quyết định trực tiếp kích thước cây tìm kiếm và hiệu quả của thuật
toán quay lui.

\section{Các nguyên tắc cơ bản}

Khi lựa chọn thứ tự tìm kiếm, ta thường dựa vào ba nguyên tắc sau.

\begin{enumerate}
  \item Phần tử tìm kiếm có miền giá trị nhỏ nên được xét trước.
  \item Sau khi một phần tử tìm kiếm được xác định, mức độ nó làm thay đổi
  miền giá trị của các phần tử khác được gọi là \term{lực ràng buộc}. Phần tử
  có lực ràng buộc lớn nên được xét trước.
  \item Xét trước những phần tử ảnh hưởng mạnh tới nghiệm giúp hàm đánh giá
  trong bước cắt tỉa chính xác hơn, từ đó làm cắt tỉa hiệu quả hơn.
\end{enumerate}

Ba nguyên tắc này không thay thế phân tích bài toán cụ thể, nhưng chúng là
điểm xuất phát tự nhiên để thu hẹp cây tìm kiếm.

\section{Ví dụ: giải mã toán tử}

Bài toán \term{giải mã toán tử} trong NOI 2000 có thể mô tả ngắn gọn như sau.
Các toán tử cổ được ký hiệu bằng các chữ cái thường từ \texttt{a} đến
\texttt{m}; chúng lần lượt tương ứng với một trong các ký hiệu hiện đại
\[
  0,1,2,3,4,5,6,7,8,9,+,*,=.
\]
Cho một nhóm đẳng thức viết bằng các toán tử cổ. Nếu tồn tại một ánh xạ làm
mọi đẳng thức đúng, cần in ra tất cả các quan hệ tương ứng chắc chắn xác định
được giữa toán tử cổ và toán tử hiện đại; nếu không tồn tại ánh xạ nào, in
\texttt{noway}.

\subsection{Ba phần tử đặc biệt nhất}

Trong bài này, ba ký hiệu đặc biệt nhất là \(=\), \(*\) và \(+\). Chúng phải
thỏa các điều kiện sau.

\begin{enumerate}
  \item Ba ký hiệu này không được xuất hiện ở đầu trái hoặc đầu phải của một
  đẳng thức.
  \item Không có hai ký hiệu nào trong ba ký hiệu này được đứng cạnh nhau.
  \item Ký hiệu \(=\) là ký hiệu đặc biệt nhất: trong mỗi đẳng thức nó phải
  xuất hiện đúng một lần.
\end{enumerate}

\subsection{Xác định thứ tự tìm kiếm}

Xét về miền giá trị, ba ký hiệu \(=\), \(+\), \(*\) có miền nhỏ nhất trong mọi
toán tử. Xét về lực ràng buộc, \(=\), \(+\), \(*\) rõ ràng mạnh hơn mười chữ
số \(0\) đến \(9\). Xét theo lợi ích đối với cắt tỉa, ba ký hiệu này cũng ảnh
hưởng lớn nhất tới nghiệm. Vì vậy chúng nên được đặt ở đầu dãy tìm kiếm.

Trong ba ký hiệu đó, \(=\) chịu nhiều hạn chế hơn và có miền giá trị nhỏ hơn,
nên phải được tìm trước. Thứ tự tốt nhất rút ra từ phân tích này là:
\[
  =,\quad +,\quad *,\quad \text{rồi đến mười chữ số.}
\]

\section{Hai cách giảm quy mô cây tìm kiếm}

Có hai hướng hiện thực cụ thể để giảm quy mô cây tìm kiếm bằng thứ tự tìm kiếm.

\begin{enumerate}
  \item \term{Tối ưu tĩnh thứ tự tìm kiếm}. Ví dụ: bài \term{Prime Square}
  trong IOI 1994 và bài \term{Fence} trong USACO Training.
  \item \term{Điều chỉnh động thứ tự tìm kiếm}. Ví dụ: bài duyệt bàn cờ và
  bài \term{Basketball Championship} trong BOI 1998.
\end{enumerate}

\subsection{Tối ưu tĩnh thứ tự tìm kiếm}

Trong một số bài toán, lực ràng buộc và miền giá trị của các phần tử tìm kiếm
không thay đổi nhiều trong quá trình tìm, hoặc sự thay đổi ấy không ảnh hưởng
đáng kể tới hiệu quả. Nếu phải phán đoán động miền giá trị và lực ràng buộc
của từng phần tử thì chi phí có thể lớn, trong khi hiệu quả tối ưu không cao.

Trong trường hợp đó, chỉ cần xác định thứ tự tìm kiếm trước khi bắt đầu, và
không cần thay đổi thứ tự trong quá trình tìm kiếm.

\subsection{Điều chỉnh động thứ tự tìm kiếm}

Đôi khi miền giá trị và lực ràng buộc của phần tử tìm kiếm biến đổi mạnh trong
quá trình tìm, và những biến đổi này trực tiếp ảnh hưởng đến quy mô cây tìm
kiếm. Khi đó ta cần điều chỉnh thứ tự tìm kiếm một cách động, tức dùng tìm
kiếm theo heuristic.

Tìm kiếm heuristic vẫn giữ các ưu điểm của quay lui: dùng ít bộ nhớ và lập
trình tương đối đơn giản. Ưu điểm lớn nhất của nó là có thể tìm được một
nghiệm khả thi trong thời gian ngắn, nên rất thích hợp cho lớp bài toán tìm
kiếm chỉ yêu cầu tìm một nghiệm.

\section{Ví dụ tối ưu tĩnh: Prime Square}

Bài \term{Prime Square} trong IOI 1994 yêu cầu điền chữ số vào một ma trận
\(5\times 5\) sao cho năm chữ số trên mỗi hàng, mỗi cột và hai đường chéo đều
tạo thành một số nguyên tố năm chữ số; chữ số đầu của mỗi số nguyên tố không
được là \(0\). Tổng các chữ số của mọi số nguyên tố này phải bằng một hằng số
cho trước. Hằng số ấy và chữ số ở góc trái trên của ma trận được cho trước.
Nếu có nhiều nghiệm thì phải tìm tất cả.

\subsection{Tính chất của phần tử tìm kiếm}

Các phần tử tìm kiếm trong bài này có những tính chất khác nhau.

\begin{enumerate}
  \item Hàng cuối và cột cuối: chữ số có thể điền chỉ thuộc tập
  \(\{1,3,7,9\}\), vì chúng là chữ số cuối của số nguyên tố năm chữ số.
  \item Hai đường chéo: chúng liên quan tới toàn bộ tập số nguyên tố năm chữ
  số có tổng chữ số bằng hằng số đã cho.
  \item Các hàng và cột còn lại: mức độ đặc biệt phụ thuộc vào số ô đã được
  xác định trên hàng hoặc cột đó.
\end{enumerate}

\subsection{Thứ tự tìm kiếm}

Dựa trên miền giá trị và lực ràng buộc, ta có thứ tự tìm kiếm sau.

\begin{enumerate}
  \item Hàng cuối và cột cuối là các phần tử tìm kiếm có miền giá trị nhỏ
  nhất. Chúng cũng ràng buộc ở mức nhất định lên mọi phần tử khác, nên cần
  đặt ở đầu dãy tìm kiếm.
  \item Hai đường chéo cũng ảnh hưởng tới mọi phần tử tìm kiếm khác. Trong
  các ô còn lại, lực ràng buộc của chúng là lớn nhất, vì vậy chúng cũng nên
  được ưu tiên.
  \item Các hàng và cột còn lại được sắp theo kích thước miền giá trị của
  chúng.
\end{enumerate}

\section{Ví dụ điều chỉnh động: Basketball Championship}

Bài \term{Basketball Championship} trong BOI 1998 có thể mô tả như sau. Có
năm đội tham gia một giải bóng rổ vòng tròn một lượt; mỗi cặp đội đấu đúng
một trận. Ta biết số trận thắng, số trận thua, tổng điểm ghi được và tổng
điểm bị ghi của mỗi đội, đồng thời biết danh sách điểm số của mọi trận đấu.
Yêu cầu là khôi phục các khả năng có thể xảy ra của giải đấu, tức bảng thành
tích điểm số của từng đội trong từng trận.

\subsection{Căn cứ để điều chỉnh động}

Trong bài này, lực ràng buộc giữa các phần tử tìm kiếm, tức các điểm số của
từng trận, rất khó xác định chính xác.

Ta định nghĩa miền giá trị của một phần tử là số cách có thể chọn điểm của
một đội trong một trận cụ thể. Vì miền giá trị của điểm số thay đổi mạnh trong
quá trình tìm, và sự thay đổi ấy trực tiếp ảnh hưởng đến quy mô cây tìm kiếm,
căn cứ chính để thay đổi thứ tự tìm kiếm là kích thước miền giá trị của từng
điểm số.

\subsection{Phán đoán động miền giá trị}

Trước hết cần tiền xử lý: liệt kê mọi trường hợp mà tổng của \(n\) lần ghi
điểm bằng \(m\). Có thể dùng bảng băm để tăng tốc tra cứu. Trong mỗi bước tìm
kiếm, dựa trên các điểm số đã điền, ta có thể xác định trong các điểm số còn
lại những giá trị nào không thể đặt vào một vị trí cụ thể.

Ngoài ra, nếu đã biết kết quả thắng thua của trận giữa \(i\) và \(j\), và đã
biết số điểm mà \(i\) ghi được trong trận đó, thì điểm số mà \(j\) ghi vào
rổ của \(i\) cũng bị hạn chế. Từ đó có thể xác định miền giá trị của từng
phần tử.

\subsection{So sánh thời gian chạy}

\begin{center}
\begin{tabular}{c|cc|cc}
\hline
Dữ liệu &
\multicolumn{2}{c|}{Tìm một nghiệm} &
\multicolumn{2}{c}{Tìm mọi nghiệm} \\
 & Tìm thường + cắt tỉa & Điều chỉnh động
 & Tìm thường + cắt tỉa & Điều chỉnh động \\
\hline
1 & 2.09 sec & 0.00 sec & 12.97 sec & 0.11 sec \\
2 & 2.69 sec & 0.00 sec & 15.22 sec & 0.00 sec \\
3 & quá thời gian & 0.16 sec & quá thời gian & 0.11 sec \\
4 & quá thời gian & 0.05 sec & quá thời gian & 0.27 sec \\
5 & quá thời gian & 0.27 sec & quá thời gian & 1.10 sec \\
\hline
\end{tabular}
\end{center}

Bảng trên cho thấy việc điều chỉnh động thứ tự tìm kiếm có thể biến một cây
tìm kiếm quá lớn thành một quá trình xử lý rất ngắn.

\section{Thứ tự tìm kiếm và tối ưu cắt tỉa}

Tối ưu cắt tỉa luôn là một bước quan trọng nhất trong tối ưu tìm kiếm. Tuy
nhiên, hiệu quả của cắt tỉa phụ thuộc vào thứ tự tìm kiếm. Chỉ khi chọn đúng
thứ tự có lợi cho cắt tỉa, các điều kiện cắt tỉa mới phát huy đầy đủ tác dụng.

\subsection{Ví dụ: Birthday Cake}

Bài \term{Birthday Cake} trong NOI 1999 yêu cầu thiết kế một chiếc bánh gồm
\(m\) tầng, có thể tích \(N\pi\). Mỗi tầng là một hình trụ; chiều cao \(H\)
và bán kính \(R\) của mỗi tầng đều nhỏ hơn tầng bên dưới nó. Với \(N\) và
\(m\) cho trước, hãy tìm diện tích mặt ngoài nhỏ nhất có thể của chiếc bánh,
không tính đáy dưới của tầng thấp nhất.

\subsection{Ảnh hưởng của thứ tự tìm kiếm tới cắt tỉa}

Vì \(R\) và \(H\) giảm dần từ dưới lên trên, thể tích mỗi tầng cũng giảm dần
từ dưới lên trên. Do đó, chọn thứ tự tìm kiếm từ dưới lên giúp phán đoán tình
trạng hiện tại có còn khả thi hay không.

Ngoài ra, diện tích mặt ngoài gồm diện tích xung quanh của từng tầng và diện
tích phần mặt trên lộ ra. Khi tầng thấp nhất đã được xác định, tổng diện tích
mặt trên lộ ra của mọi tầng cũng được xác định; phần còn phải xét chỉ là diện
tích xung quanh của từng tầng. Vì vậy thứ tự tìm kiếm này cũng giúp phán đoán
liệu còn khả năng xuất hiện nghiệm tốt hơn nghiệm tốt nhất đã biết hay không.

\section{Kết luận}

Trong giải bài thực tế và trong ứng dụng, bài toán tìm kiếm hầu như luôn gắn
với tối ưu, mà bước đầu tiên của tối ưu là chọn đúng thứ tự tìm kiếm. Từ các
phân tích trên có thể thấy lựa chọn thứ tự tìm kiếm có vai trò quan trọng đối
với hiệu quả của thuật toán tìm kiếm.

Tuy nhiên, việc chọn thứ tự tìm kiếm không thể chỉ dựa vào vài phương pháp và
tư tưởng cơ bản. Nó đòi hỏi phân tích bài toán một cách toàn diện và thấu
đáo, đào sâu bản chất của bài toán, từ đó tìm ra điểm đột phá của lời giải.

\end{document}
